\section{Task and Data}

We use COLIEE 2026 Task~1 as the main benchmark and COLIEE 2025 as auxiliary data for methodological exploration and failure analysis. The goal of Task~1 is not to retrieve generally similar judgments, but to find the noticed cases for each query case. If a case was referenced by the query case, we treat it as a noticed case for that query. Because the organizers redact citation information from all case texts, the system must infer from the facts, reasoning, and legal semantics which prior cases most likely support the judgment.

\subsection{Data Sources and Splits}

The data come from the Federal Court of Canada. In the training data, the organizers provide query cases, the full candidate pool, and the gold noticed cases for each query. In the test data, only query cases are released, and the system must automatically predict the noticed cases. In other words, the whole pipeline from retrieval to final output must be fully automatic.

Table~\ref{tab:dataset-stats} summarizes the query counts, year-specific candidate pool sizes, and average numbers of noticed cases per query. COLIEE 2025 provides labeled train and validation splits, while for COLIEE 2026 we use the released labeled training set together with the unlabeled official test queries. In our current setup, the COLIEE 2026 case collection contains 7,708 documents. We further split the official labels into 1,600 training queries and 401 validation queries. Model training and hyperparameter selection are performed on this train/validation split. Our main analysis reports top-5 micro-averaged F1, precision, and recall on both the 2026 validation and 2026 training queries. COLIEE 2025 is mainly used to document training failure, the effect of dynamic negatives, and the early gain from scope filtering.

\begin{table*}[t]
\caption{Dataset statistics of COLIEE Task~1. Candidate counts are measured from the year-specific cleaned full-document corpus.}
\label{tab:dataset-stats}
\centering
\small
\setlength{\tabcolsep}{5pt}
\begin{tabular}{lccccc}
\toprule
 & \multicolumn{2}{c}{COLIEE 2025} & \multicolumn{3}{c}{COLIEE 2026} \\
\cmidrule(lr){2-3} \cmidrule(lr){4-6}
Statistic & Train & Valid & Train & Valid & Test \\
\midrule
\# of query cases & 1,341 & 336 & 1,600 & 401 & 400 \\
\# of candidate cases & 7,349 & 7,349 & 7,708 & 7,708 & 7,708 \\
Avg. \# of noticed cases per query & 4.07 & 4.24 & 4.12 & 4.13 & 4.38 \\
\bottomrule
\end{tabular}
\end{table*}

\subsection{Full-Document and Citation-Context Query Views}

We compare two query views. The first is the full-document view, which uses the full cleaned judgment as both query and candidate. The second is the citation-context view, which follows THUIR's reference sentence extraction idea and keeps local context around the citation-redaction placeholders; when the extracted context is too short, it is expanded with the beginning of the cleaned full judgment. In our study, this is a query hypothesis that must be re-tested rather than accepted by default. For fair ablation, candidate documents are always encoded from the full-document view; only the query view changes.

The full-document view keeps the main judgment body while removing suppression placeholders such as \texttt{FRAGMENT\_SUPPRESSED}, editorial markers such as \texttt{[translation]}, \texttt{[sic]}, and \texttt{[End of document]}, and passages detected as French. If the cleaned text still lacks a \texttt{Summary:} section, an existing summary block is prepended only when that summary already appears in the raw case file. The citation-context view first extracts sentences around citation-redaction placeholders, then applies nearly the same cleanup steps, and finally expands very short queries with the beginning of the cleaned full judgment. Thus, it is not hard-capped at 512 words; 512 English word tokens only serves as the threshold for fallback expansion. On COLIEE 2026, the full-document view averages 3,965.52 words, whereas the citation-context view averages 882.54 words with median 484. Therefore, the second view is not a pure citation fragment but a citation-context query view with fallback expansion from the cleaned full judgment.

\subsection{Long-Document Characteristics and Continued Pretraining Data}

As shown in Table~\ref{tab:length-stats}, COLIEE 2026 cleaned full-document cases are clearly long, and 94.80\% of them are shorter than 12,288 tokens. This makes three 4096-token chunks a practical reranking approximation for most cases.

Besides the official COLIEE data, we use the Common Pile CaseLaw Access Project subset from Hugging Face\footnote{\url{https://huggingface.co/datasets/common-pile/caselaw_access_project_filtered}} for continued pretraining. It contains 6,735,525 U.S. case documents spanning roughly 365 years. This corpus does not provide direct COLIEE supervision; instead, it is used for DAPT before contrastive fine-tuning. Under a single RTX 4080 SUPER 16GB and our time budget, this DAPT stage runs for about 60 hours and still covers only 0.384 epoch, so we treat it as a useful but clearly under-trained adaptation step.
